% =============================================================================
%  Lab 5 - Arithmetic Logic Unit
% =============================================================================

\chapter{Lab 5 --- Arithmetic Logic Unit}

\section{Objective}
% Two or three sentences, your own words.

\section{Specification}

\begin{table}[htbp]
\centering
\caption{ALU operations. One worked example per operation --- fill in the last
column from your own simulation.}
\label{tab:lab05-ops}
\begin{tabular}{lll@{\hskip 10pt}l}
\toprule
code & operation & \texttt{Rout} & worked example \\
\midrule
\texttt{ALU\_FUNC\_ADD}  & add                & \texttt{Ain + Bin}   & \\
\texttt{ALU\_FUNC\_SUB}  & subtract           & \texttt{Ain - Bin}   & \\
\texttt{ALU\_FUNC\_NEG}  & negate             & \texttt{0 - Bin}     & \\
\texttt{ALU\_FUNC\_AND}  & bitwise and        & \texttt{Ain and Bin} & \\
\texttt{ALU\_FUNC\_NAND} & bitwise nand       &                      & \\
\texttt{ALU\_FUNC\_OR}   & bitwise or         &                      & \\
\texttt{ALU\_FUNC\_NOR}  & bitwise nor        &                      & \\
\texttt{ALU\_FUNC\_XOR}  & bitwise xor        &                      & \\
\texttt{ALU\_FUNC\_XNOR} & bitwise xnor       &                      & \\
\texttt{ALU\_FUNC\_NOT}  & invert \texttt{Bin} &                     & \\
\texttt{ALU\_FUNC\_SLL}  & shift left logical &                      & \\
\texttt{ALU\_FUNC\_SLA}  & shift left arith.  & same as \texttt{SLL} & \\
\texttt{ALU\_FUNC\_SRL}  & shift right logical &                     & \\
\texttt{ALU\_FUNC\_SRA}  & shift right arith. &                      & \\
\texttt{ALU\_FUNC\_SOLS} & set on less, signed &                     & \\
\texttt{ALU\_FUNC\_SOLU} & set on less, unsigned &                   & \\
\bottomrule
\end{tabular}
\end{table}

\subsection*{Flag equations}

\begin{align*}
\text{Zero} &= \overline{R_{31} + R_{30} + \cdots + R_{0}} \\
\text{Neg}  &= R_{31}
\end{align*}

% ONE SENTENCE, required: why does Neg need no computation at all?

\subsection*{Set on less than}

\[
\text{SOLS} = R_{31} \oplus V \qquad \text{where } R = \texttt{Ain} - \texttt{Bin}
\]

% REQUIRED, two or three sentences: the naive answer is "subtract and look at
% the sign bit". Explain exactly when that fails and why the XOR with the
% overflow flag repairs it.

\section{Design}
% Answer all three:
%   (a) Your ALU is a 'case' statement. Explain why that IS a multiplexer, and
%       point at your RTL Viewer figure as evidence.
%   (b) Why are the flags computed AFTER the case, on the selected result?
%       Give the specific operation and operands that would break if you took
%       the adder's flags instead.
%   (c) SRA and SRL use the same VHDL operator. What makes them behave
%       differently, and why does that make the bug so easy to miss?

\section{Implementation}

\vhdlfile{code/lab05/ALU.vhd}{\texttt{ALU} --- sixteen operations and a multiplexer}

\section{Verification}
% Vectors, then the coverage argument.
%
% The argument for ARITHMETIC is different from Lab 3's. There is no selector to
% enumerate. State it as BOUNDARIES:
%     zero, one, 0x7FFFFFFF, 0x80000000, 0xFFFFFFFF
% and say why the middle of the range is uninteresting.

\begin{table}[htbp]
\centering
\caption{Test vectors. The last column is the point of the table.}
\label{tab:lab05-vectors}
\begin{tabular}{clllp{5cm}}
\toprule
\# & operation & operands & expected & which wrong implementation this catches \\
\midrule
1 & & & & \\
2 & & & & \\
3 & & & & \\
\bottomrule
\end{tabular}
\end{table}

\section{Results}
% Evidence: E1, E2, E4, E6, E7. Two of these are new this week.

\figplaceholder{the RTL Viewer for \texttt{ALU}, with the multiplexer visible}
% Required sentence: identify the multiplexer the synthesiser drew from your
% case statement.

\figplaceholder{the waveform showing \texttt{SRA} and \texttt{SRL} applied to the
same negative value}
% Required sentence: state both results and say which is the correct arithmetic
% shift, with the decimal values.

% E4: ALM count, and the register count. The register count MUST be zero -
%     say why a non-zero count would mean the design is wrong, not just costly.
% E6: Fmax on the Slow 1100mV 85C model, and the critical path. Say which
%     operation the path runs through and why that is unsurprising. Hint: the
%     logic operations are one gate deep; something here is 32 gates deep.

\section{Discussion}
% REQUIRED. One bug, how you found it, what it cost.
%
% If you are short of a bug, a good honest topic: did your first SRA use
% unsigned()? How did you find out - and would your original test vectors have
% caught it?
